\begin{figure}[t]
\centering
\begin{tikzpicture}[>=stealth,font=\footnotesize,
box/.style={draw=sgblue!50,fill=sgblue!3,align=center,text width=7.3cm,minimum height=.55cm,inner sep=3pt}]
\node[box] (h0) at (0,0.0) {$N=4$: 16 public examples};
\node[box] (h1) at (0,-0.85) {Consistent pool: 10--14 functions};
\node[box] (h2) at (0,-1.7) {Exactly 2 actions: target shares $q$ and $1-q$};
\draw[->] (h0.south)--(h1.north);\draw[->] (h1.south)--(h2.north);
\node at (0,-2.35) {$q=.8,.7,.6,.5$; three rounds each};
\node[box,fill=sgorange!5,draw=sgorange!60] at (0,-3.25) {Fewer participants: shorter prefix\\$N: 4\to3\to2\to1$; examples: $16\to12\to8\to4$};
\node[font=\scriptsize] at (0,-4.05) {Pool: $\sim10\to50$--$70\to200$--$370\to2{,}227$};
\end{tikzpicture}
\caption{Public-evidence mechanism. The truth is placed on the majority side with probability $q$ before redraws. Splits are nominal targets; finite-pool diagnostics are reported in the text. Fewer participants retain a prefix, leaving more hypotheses. Counts are supplied design diagnostics, not new estimates.}
\label{fig:arena-hints}
\end{figure}
